\section{Additional Interpretable Audio Phenomena: Prospective Extension}
\label{sec:additional-phenomena}

\noindent\textbf{Dated evidence boundary (September 7, 2026).}
This section extends the magnitude-spectral (S), attack-amplitude dynamics (D),
rhythm (R), and section-length (P) families with explicitly defined candidate
phenomena. It is exploratory development work on previously assembled sources,
not a new untouched external attribution test. Historical locked, stress and
pilot labels do not enter fitting or feature selection. The September 5 report
and source audio remain unchanged. Successful measurement of a phenomenon and
successful Human/AI classification are different hypotheses.

\subsection{Literature-to-measurement translation}

The candidate registry comprises phase evolution/group delay (F), tonal/chroma
paths (H), long-range content recurrence (M), pitch microstructure (V), audible
breath organization (B), articulation/phone--note coordination (A), and
source-conditioned timbral continuity (T). A new representation or a new name
for spectral flux is not automatically an independent phenomenon.

Direct music-detection evidence motivates F: Afchar et al.,
\emph{AI-Generated Music Detection and its Challenges} (ICASSP 2025), demonstrate
detectable information in a phase representation under controlled reconstruction.
This does not validate a universal hand-designed phase statistic.
\footnote{\url{https://arxiv.org/abs/2501.10111}}
H and M are supported by symbolic evaluation and content-recurrence analysis,
including Wu and Yang, \emph{The Jazz Transformer on the Front Line} (ISMIR 2020),
and Kim and Go, \emph{Segment Transformer} (APSIPA ASC 2025). Their tasks and
representations differ from this mixed-waveform experiment.
\footnote{\url{https://archives.ismir.net/ismir2020/paper/000339.pdf};
\url{https://www.apsipa.org/proceedings/2025/papers/APSIPA2025_P179.pdf}}
Pitch trajectories, breath events and timbral consistency have measurement or
perceptual motivation; this is not in itself evidence that contemporary AI
music lacks them. Primary papers, proposed-versus-published metrics, counterexamples
and publication status are retained in
\filepath{../../interpretable_phenomena_literature_20260906/RESEARCH_NOTE_EN.md}.

\begin{longtable}{>{\raggedright\arraybackslash}p{29mm}>{\raggedright\arraybackslash}p{61mm}>{\raggedright\arraybackslash}p{59mm}}
\caption{Additional phenomena: measurement targets and the actual evidence level. Proposed metrics are hypotheses unless explicitly attributed to a published measure.}\\
\toprule Phenomenon & Interpretable measurement target & Publication support and boundary \\
\midrule\endfirsthead
\toprule Phenomenon & Interpretable measurement target & Publication support and boundary \\
\midrule\endhead
F: phase evolution & Frame-to-frame circular phase residual, group-delay dispersion, attack--decay contrast & Afchar et al., ICASSP 2025: direct music reconstruction-detection representation evidence, not validation of these scalar statistics. \\
V: pitch microstructure & Note-relative cents, vibrato rate/extent, approach overshoot, glide shape & Saitou et al., Speech Communication 2005; Abe\ss er et al., ISMIR 2015: perception and note-level performance measurement, not AI attribution. \\
B: breath organization & Audible inhale-event rate and duration, intervals and alignment to vocal phrase boundaries & Layton et al., DTRAP 2025; Mostaani and Magimai-Doss, Interspeech 2022: fake-speech evidence requiring separate singing validation. \\
A: articulation coordination & Phone durations conditioned on notes, vowel/consonant transitions and melisma & Wang et al., Interspeech 2023: learned duration/pronunciation features for fake speech; phone--note coordination is our extension hypothesis. \\
H: tonal/harmonic paths & Pitch-class entropy, chord-trigram irregularity, local tonal-state transitions & Wu and Yang, ISMIR 2020: interpretable symbolic jazz-generation evaluation, not a universal harmony-error rule. \\
M: motif recurrence & Off-diagonal self-similarity, recurrence lag and returning-content variation & Kim and Go, APSIPA ASC 2025: direct music-detection evidence using learned segment embeddings and self-similarity; hand-designed recurrence still needs validation. \\
T: timbral identity & Register/intensity-conditioned similarity within a source or recurring note & Nercessian et al., ISMIR 2024: cross-note instrument sample-bank consistency, not whole-song source tracking or a validated AI detector. \\
\bottomrule
\end{longtable}

The V publications explicitly support pitch-performance analysis rather than
authorship detection. Saitou et al. studied overshoot, vibrato, preparation and
fine fluctuation in singing synthesis; Abe\ss er et al. studied score-informed
jazz intonation and pitch modulation.
\footnote{\url{https://doi.org/10.1016/j.specom.2005.01.010};
\url{https://archives.ismir.net/ismir2015/paper/000135.pdf}}
The final breath article is \emph{Every Breath You Don't Take: Deepfake Speech
Detection Using Breath}, Digital Threats: Research and Practice 6(3), article
13 (2025), not a singing study. Mostaani and Magimai-Doss also distinguish
the limitations of breath information for converted natural speech.
\footnote{\url{https://doi.org/10.1145/3754456};
\url{https://doi.org/10.21437/Interspeech.2022-10271}}
Wang et al.'s pronunciation representation is learned; it does not validate
consonant-release timing as a hand-designed singing detector.
\footnote{\url{https://doi.org/10.21437/Interspeech.2023-1254}}
Nercessian et al.'s timbral-consistency metric compares CLAP affinity structure
across pitch/velocity sample sets. Moving from a sample bank to a temporal
source-identity trajectory is a new hypothesis.
\footnote{\url{https://ismir2024program.ismir.net/poster_22.html}}
None of these papers establishes a universal direction such as ``more regular
means AI.'' Human pitch correction, loop composition, de-breathing and source
separation are important counterexamples. Fixed spectral grids overlap S and
are not counted again as an independent phenomenon.

\subsection{Exact audio inputs and independent measurement validation}

F/H use the original mix, not the Demucs vocal. Inputs preserve the exact
historical crop offset, use arithmetic channel averaging, remove DC and resample
with a fixed polyphase Kaiser-5 filter to 16\,kHz. Native rate must be at least
16\,kHz. Source SHA-256, analyzed-waveform SHA-256, crop frames, code/configuration
hashes and runtime versions are retained per recording. No short input is padded.
The previous magnitude-domain separator correction is not a phase or F0 correction.

\begin{table}[htbp]
\centering\small
\caption{Completed extraction accounting. Observable means an eligible signal
measurement, not a successful classifier prediction.}
\label{tab:phenomena-extraction}
\begin{tabular}{lrrrr}
\toprule Fixed view & Records & F observable & H observable & M observable \\
\midrule
10 seconds & 10,141 & 10,041 & 9,323 & 0 \\
30 seconds & 4,497 & 4,496 & 4,425 & 0 \\
60 seconds & 2,092 & 2,092 & 2,076 & 2,055 \\
\bottomrule
\end{tabular}
\end{table}

All three extractions have complete per-item accounting and no processing failure.
F non-observability consists of low energy or insufficient phase coverage; H
requires informative pitch-class structure. These cases remain missing, not
zero-valued evidence. M requires at least 45 seconds and is not evaluated on
the short views. A separate exact-60-second cohort contains 2,092 physically
verified crops (1,546 Human, 546 AI). Its only development AI source is Suno;
the 50 DiffRhythm recordings remain pilot data. Therefore that cohort cannot
support unseen-generator generalization of long recurrence. Longer context and
changed source composition are not an isolated duration intervention.

F contains 15 statistics of weighted circular phase residuals and robust group
delay, summarized over all, attack, sustain and decay regions and selected
decay-minus-sustain contrasts. At 16\,kHz the periodic-Hann STFT has 1,024 samples
and a 256-sample hop. Energy masks exclude unreliable near-zero bins. With the
forward-FFT convention, the residual is
\begin{equation}
 r_{k,t}=\operatorname{wrap}\left(\phi_{k,t+1}-\phi_{k,t}
                 -2\pi kH/N\right),\qquad
 v=1-\left|\frac{\sum w_{k,t}e^{\mathrm{i}r_{k,t}}}{\sum w_{k,t}}\right|.
\end{equation}
Phase/group-delay features share one family; they are not treated as 15
independent biological hypotheses. Attack/decay regions are energy-direction
proxies, not manually annotated musical articulations.

External F controls use 144 independent MUSDB tracks and 1,632 matched pairs,
including phase interventions, gain/polarity, time shifts, MP3 and reference
versus separated vocals. Gain/polarity invariance is numerically strong, but
shift, codec and separation changes are detectable. Preserving global Fourier
magnitude in a phase intervention does not preserve every local STFT magnitude,
transient or envelope. Consequently the result establishes sensitivity and
confounding risks, not a pure causal generator signature.

H contains six pitch-class entropy and chroma-path descriptors. A 4,096-point
STFT with a 1,024-sample hop feeds one-second normalized chroma summaries;
high-entropy/uninformative frames are excluded. M compares four-second patterns
at specified lags, allowing global pitch-class transposition. It is a
content-recurrence descriptor, not motif transcription or full-song form.
External real-audio controls materialize 36 continuous 64-second stimuli in
12 trials from 24 MUSDB recordings. Known repeated-block similarity exceeds
matched shuffled-block similarity in 12/12 trials (median difference 0.1151);
changing versus stationary H paths differ positively in 11/11 observable trials
(one unavailable; median difference 0.2680). These controlled mosaics do not
replace natural annotated chord or motif validation.

\subsection{External isolated-vocal pitch tracking}

An independently pinned MIR-1K subset provides 100 clips from 19 singers,
824.99 seconds and 41,102 reference frames. All selected WAV and pitch-label
files were hash-verified. A misleading mirror README describing a different
corpus was retained but not used: MIR-1K reference timestamps start at 20\,ms
and repeat every 20\,ms. The original right channel contains isolated singing.
The first baseline uses pYIN at 16\,kHz, a 1,024-sample frame and a 160-sample
hop, with fixed C2--C6 range. Reference-grid alignment never bridges unvoiced
gaps. No AI/Human label is used for configuration or scoring.

\begin{table}[htbp]
\centering\small
\caption{External pYIN benchmark; all 100 clips completed. These are F0-tracking
metrics, not AI-music detection accuracy.}
\label{tab:external-pyin-phenomena}
\begin{tabular}{lrr}
\toprule Metric & Value & Numerator / denominator \\
\midrule
Raw pitch accuracy, 50 cents & 0.966299 & 28,357 / 29,346 \\
Raw chroma accuracy, 50 cents & 0.966810 & 28,372 / 29,346 \\
Voicing precision & 0.881141 & 29,164 / 33,098 \\
Voicing recall & 0.993798 & 29,164 / 29,346 \\
Voicing F1 & 0.934085 & combined precision and recall \\
Overall accuracy & 0.880225 & 36,179 / 41,102 \\
\bottomrule
\end{tabular}
\end{table}

The predeclared limited gate (pitch accuracy $\geq0.70$, voicing F1 $\geq0.80$,
complete verified clips) passes. This only supports coarse F0 trajectories on
comparable isolated vocals. Vibrato rate/extent, onset overshoot, note segmentation
and Demucs robustness still need separate validation. In particular, the
false-positive voicing count matters for apparent small-scale pitch modulation.
The mirrored audio's redistribution license remains unresolved; it is not
included in a public dataset release.

The next, separately frozen V assay fits 4--8\,Hz sinusoidal modulation plus
a linear trajectory in contiguous one-second voiced windows. It rejects
large pitch steps, weak modulation and poor sinusoidal fit. Eleven analytic
test cases pass, but the real measurement gate fails: reference contours have
only three observable windows in two clips; estimated contours have four
windows in two clips; only one clip is paired. Observability precision and
recall are both 0.50, below 0.70, and the minimum 25 paired clips is not met.
The single paired error cannot establish reliability. V is therefore not
admitted to the classifier, and no threshold is loosened after seeing this
benchmark. Good coarse F0 accuracy does not validate small-scale performance
features.

\subsection{Matched classification design and additional information}

The 30-second analysis evaluates all 63 nonempty S/D/R/P/F/H combinations at
25, 50, 100, 200 and all global groups per training source. Only the 3,317
development recordings enter fitting. Group identities are blocked globally,
including shared prompts and groups crossing sources or classes. Source holdout
and ordinary recording-group cross-validation are reported separately. A fixed
ridge model, training-only imputation/scaling and a fixed 0.5 decision threshold
are used throughout. The ridge score is not a calibrated probability.

Define $J$ as the equal mean of Human-source-heldout and AI-generator-heldout
macro AUC. Source-pair results are aggregated within held-out source before
macro averaging. The table below uses identical realized rows, folds and group-cap
policy for every matched increment; subtracting a September 5 result obtained
under a different evaluation configuration would not be a valid increment.
\begin{table}[htbp]
\centering\small
\caption{Selected exact-30-second results at the all-groups training cap. AUC columns are source-macro; BA uses the fixed 0.5 threshold. These are development, not untouched test results.}
\label{tab:fh30-main}
\begin{tabular}{lrrrr}
\toprule Combination & Human-held AUC & AI-held AUC & $J$ & Mean BA \\
\midrule
S & 0.7216 & 0.7056 & 0.7136 & 0.6590 \\
D & 0.6673 & 0.6701 & 0.6687 & 0.6408 \\
R & 0.6571 & 0.6583 & 0.6577 & 0.6091 \\
P & 0.5554 & 0.5274 & 0.5414 & 0.5115 \\
F & 0.6144 & 0.6356 & 0.6250 & 0.6025 \\
H & 0.5788 & 0.5462 & 0.5625 & 0.5407 \\
F+H & 0.6404 & 0.6661 & 0.6533 & 0.6239 \\
S+D+R & 0.7623 & 0.7595 & 0.7609 & 0.6986 \\
S+D+R+F & 0.7654 & 0.7710 & 0.7682 & 0.7010 \\
S+D+R+H & 0.7559 & 0.7557 & 0.7558 & 0.6909 \\
S+D+R+F+H & 0.7647 & 0.7714 & 0.7680 & 0.6994 \\
S+D+R+P+F+H & 0.7687 & 0.7595 & 0.7641 & 0.6972 \\
\bottomrule
\end{tabular}
\end{table}

\begin{table}[htbp]
\centering\small
\caption{Predefined paired increments over S+D+R at all groups: 1,000 global-group bootstrap replicates, conditional on frozen models and the observed sources. Intervals are unadjusted percentile intervals.}
\label{tab:fh30-group-ci}
\begin{tabular}{lrrrr}
\toprule Added & $\Delta J$ & 95\% interval & $\Delta$ Mean BA & 95\% interval \\
\midrule
F & +0.0073 & [-0.0046, +0.0190] & +0.0024 & [-0.0100, +0.0152] \\
H & -0.0051 & [-0.0095, -0.0006] & -0.0078 & [-0.0156, -0.0017] \\
F+H & +0.0071 & [-0.0050, +0.0203] & +0.0008 & [-0.0122, +0.0133] \\
\bottomrule
\end{tabular}
\end{table}

Neither F nor F+H has a positive interval bounded away from zero for the equal-direction criterion. H has a negative group-conditional increment here; this is not a universal statement about harmony in generated music. Five replicates contain 40 undefined source-pair cells in total; every macro remains defined using the frozen evaluator\textquotesingle s NaN-skipping rule. All paired point estimates were reproduced to $1.11\times10^{-16}$ or better.

\begin{table}[htbp]
\centering\small
\caption{Additional post-hoc empirical source-composition sensitivity: 10,000 draws of six Human and four AI source labels. The same crossed-source weights apply to both evaluation directions and all comparisons. These ranges do not combine group uncertainty and are not future-source confidence guarantees.}
\label{tab:fh30-source-sensitivity}
\begin{tabular}{lrr}
\toprule Added & $\Delta J$ empirical 95\% range & $\Delta$ Mean BA empirical 95\% range \\
\midrule
F & [-0.0459, +0.0567] & [-0.0410, +0.0466] \\
H & [-0.0241, +0.0083] & [-0.0170, +0.0021] \\
F+H & [-0.0470, +0.0582] & [-0.0416, +0.0463] \\
\bottomrule
\end{tabular}
\end{table}

All source-reweighted ranges include zero. This sensitivity was added after seeing point estimates and is explicitly post-hoc. It reinforces the source-composition limitation, not a claim of an iid population sample. Models are never refitted and threshold 0.5 remains fixed.

\begingroup\footnotesize\setlength{\tabcolsep}{4pt}
\begin{longtable}{lrrrrr}
\caption{Complete 63-combination by five-quantity development lattice. Each cell is $J$/Mean BA. Quantities cap training global groups per source, not total songs. All rows share the same 3,317-recording eligible cohort; realized training sizes vary by fold and source availability.}\label{tab:fh30-full-lattice}\\
\toprule Combination & 25 & 50 & 100 & 200 & All \\
\midrule\endfirsthead
\toprule Combination & 25 & 50 & 100 & 200 & All \\
\midrule\endhead
\midrule\multicolumn{6}{r}{Continued on next page}\\\endfoot
\bottomrule\endlastfoot
S & 0.7054/0.6514 & 0.7041/0.6550 & 0.7142/0.6591 & 0.7101/0.6543 & 0.7136/0.6590 \\
D & 0.6627/0.6283 & 0.6666/0.6315 & 0.6682/0.6389 & 0.6672/0.6375 & 0.6687/0.6408 \\
R & 0.6566/0.6067 & 0.6553/0.6104 & 0.6588/0.6127 & 0.6588/0.6125 & 0.6577/0.6091 \\
P & 0.5388/0.5281 & 0.5332/0.5136 & 0.5200/0.5017 & 0.5315/0.5054 & 0.5414/0.5115 \\
F & 0.6277/0.6102 & 0.6273/0.6065 & 0.6226/0.6044 & 0.6268/0.6042 & 0.6250/0.6025 \\
H & 0.5663/0.5482 & 0.5726/0.5484 & 0.5656/0.5399 & 0.5600/0.5356 & 0.5625/0.5407 \\
S+D & 0.7177/0.6650 & 0.7169/0.6655 & 0.7255/0.6670 & 0.7221/0.6625 & 0.7286/0.6668 \\
S+R & 0.7320/0.6646 & 0.7311/0.6712 & 0.7433/0.6786 & 0.7436/0.6788 & 0.7462/0.6806 \\
S+P & 0.7014/0.6527 & 0.6945/0.6461 & 0.7049/0.6518 & 0.7085/0.6546 & 0.7148/0.6586 \\
S+F & 0.7296/0.6735 & 0.7222/0.6622 & 0.7283/0.6675 & 0.7316/0.6705 & 0.7347/0.6711 \\
S+H & 0.6914/0.6440 & 0.6955/0.6491 & 0.7081/0.6532 & 0.7039/0.6500 & 0.7084/0.6543 \\
D+R & 0.7138/0.6523 & 0.7200/0.6588 & 0.7268/0.6576 & 0.7250/0.6586 & 0.7251/0.6566 \\
D+P & 0.6396/0.6141 & 0.6489/0.6159 & 0.6513/0.6223 & 0.6648/0.6350 & 0.6697/0.6405 \\
D+F & 0.6741/0.6390 & 0.6821/0.6463 & 0.6803/0.6464 & 0.6819/0.6491 & 0.6835/0.6529 \\
D+H & 0.6321/0.5972 & 0.6483/0.6040 & 0.6467/0.5988 & 0.6416/0.5954 & 0.6421/0.5978 \\
R+P & 0.6527/0.6016 & 0.6563/0.6139 & 0.6532/0.6098 & 0.6592/0.6164 & 0.6619/0.6170 \\
R+F & 0.6477/0.6210 & 0.6500/0.6199 & 0.6467/0.6214 & 0.6504/0.6212 & 0.6514/0.6229 \\
R+H & 0.6497/0.6023 & 0.6614/0.6072 & 0.6689/0.6168 & 0.6686/0.6133 & 0.6696/0.6130 \\
P+F & 0.6172/0.5936 & 0.6260/0.6034 & 0.6183/0.5991 & 0.6274/0.6055 & 0.6304/0.6065 \\
P+H & 0.5758/0.5541 & 0.5799/0.5546 & 0.5676/0.5423 & 0.5713/0.5450 & 0.5812/0.5529 \\
F+H & 0.6400/0.6105 & 0.6465/0.6202 & 0.6485/0.6197 & 0.6545/0.6235 & 0.6533/0.6239 \\
S+D+R & 0.7426/0.6777 & 0.7436/0.6813 & 0.7554/0.6905 & 0.7556/0.6924 & 0.7609/0.6986 \\
S+D+P & 0.7084/0.6617 & 0.7038/0.6529 & 0.7146/0.6591 & 0.7194/0.6612 & 0.7286/0.6665 \\
S+D+F & 0.7410/0.6799 & 0.7368/0.6738 & 0.7448/0.6785 & 0.7483/0.6804 & 0.7551/0.6924 \\
S+D+H & 0.7026/0.6517 & 0.7070/0.6581 & 0.7189/0.6623 & 0.7148/0.6578 & 0.7218/0.6610 \\
S+R+P & 0.7219/0.6643 & 0.7204/0.6591 & 0.7300/0.6660 & 0.7384/0.6725 & 0.7433/0.6746 \\
S+R+F & 0.7412/0.6810 & 0.7364/0.6752 & 0.7422/0.6790 & 0.7467/0.6797 & 0.7496/0.6834 \\
S+R+H & 0.7197/0.6630 & 0.7244/0.6642 & 0.7389/0.6737 & 0.7402/0.6775 & 0.7443/0.6784 \\
S+P+F & 0.7200/0.6594 & 0.7078/0.6511 & 0.7169/0.6586 & 0.7262/0.6650 & 0.7335/0.6679 \\
S+P+H & 0.6892/0.6428 & 0.6859/0.6400 & 0.6991/0.6441 & 0.7018/0.6460 & 0.7087/0.6505 \\
S+F+H & 0.7145/0.6527 & 0.7145/0.6518 & 0.7267/0.6645 & 0.7342/0.6708 & 0.7383/0.6709 \\
D+R+P & 0.6938/0.6499 & 0.7076/0.6544 & 0.7106/0.6514 & 0.7186/0.6582 & 0.7213/0.6586 \\
D+R+F & 0.6892/0.6467 & 0.6994/0.6505 & 0.7001/0.6515 & 0.7020/0.6473 & 0.7054/0.6544 \\
D+R+H & 0.6891/0.6305 & 0.7068/0.6483 & 0.7161/0.6505 & 0.7128/0.6469 & 0.7143/0.6475 \\
D+P+F & 0.6573/0.6224 & 0.6700/0.6321 & 0.6693/0.6391 & 0.6771/0.6420 & 0.6831/0.6493 \\
D+P+H & 0.6206/0.5951 & 0.6368/0.6022 & 0.6358/0.5955 & 0.6419/0.6006 & 0.6479/0.6028 \\
D+F+H & 0.6683/0.6355 & 0.6824/0.6468 & 0.6852/0.6461 & 0.6880/0.6474 & 0.6910/0.6503 \\
R+P+F & 0.6354/0.6090 & 0.6464/0.6143 & 0.6383/0.6136 & 0.6472/0.6179 & 0.6526/0.6208 \\
R+P+H & 0.6463/0.6012 & 0.6605/0.6089 & 0.6605/0.6079 & 0.6676/0.6147 & 0.6718/0.6184 \\
R+F+H & 0.6604/0.6205 & 0.6666/0.6294 & 0.6714/0.6351 & 0.6784/0.6366 & 0.6804/0.6397 \\
P+F+H & 0.6297/0.5985 & 0.6418/0.6084 & 0.6427/0.6118 & 0.6538/0.6196 & 0.6572/0.6203 \\
S+D+R+P & 0.7293/0.6708 & 0.7309/0.6666 & 0.7419/0.6771 & 0.7504/0.6826 & 0.7572/0.6901 \\
S+D+R+F & 0.7523/0.6941 & 0.7506/0.6827 & 0.7576/0.6888 & 0.7624/0.6932 & 0.7682/0.7010 \\
S+D+R+H & 0.7281/0.6703 & 0.7342/0.6764 & 0.7495/0.6829 & 0.7496/0.6838 & 0.7558/0.6909 \\
S+D+P+F & 0.7291/0.6702 & 0.7202/0.6605 & 0.7310/0.6683 & 0.7414/0.6779 & 0.7516/0.6892 \\
S+D+P+H & 0.6949/0.6429 & 0.6939/0.6454 & 0.7085/0.6549 & 0.7119/0.6537 & 0.7214/0.6601 \\
S+D+F+H & 0.7255/0.6640 & 0.7281/0.6675 & 0.7399/0.6781 & 0.7463/0.6766 & 0.7534/0.6888 \\
S+R+P+F & 0.7296/0.6679 & 0.7230/0.6628 & 0.7295/0.6691 & 0.7402/0.6745 & 0.7465/0.6802 \\
S+R+P+H & 0.7123/0.6588 & 0.7142/0.6533 & 0.7260/0.6628 & 0.7344/0.6742 & 0.7410/0.6773 \\
S+R+F+H & 0.7291/0.6670 & 0.7302/0.6628 & 0.7418/0.6818 & 0.7506/0.6849 & 0.7548/0.6909 \\
S+P+F+H & 0.7056/0.6454 & 0.7005/0.6395 & 0.7153/0.6600 & 0.7292/0.6641 & 0.7372/0.6706 \\
D+R+P+F & 0.6716/0.6342 & 0.6881/0.6388 & 0.6866/0.6431 & 0.6948/0.6437 & 0.7023/0.6511 \\
D+R+P+H & 0.6744/0.6268 & 0.6952/0.6416 & 0.7014/0.6365 & 0.7066/0.6425 & 0.7112/0.6476 \\
D+R+F+H & 0.6860/0.6414 & 0.6988/0.6470 & 0.7052/0.6520 & 0.7091/0.6548 & 0.7130/0.6615 \\
D+P+F+H & 0.6537/0.6196 & 0.6707/0.6299 & 0.6749/0.6347 & 0.6844/0.6384 & 0.6906/0.6454 \\
R+P+F+H & 0.6493/0.6107 & 0.6604/0.6230 & 0.6625/0.6280 & 0.6749/0.6336 & 0.6810/0.6382 \\
S+D+R+P+F & 0.7393/0.6752 & 0.7356/0.6692 & 0.7436/0.6763 & 0.7545/0.6882 & 0.7640/0.6962 \\
S+D+R+P+H & 0.7164/0.6606 & 0.7220/0.6602 & 0.7369/0.6705 & 0.7438/0.6762 & 0.7523/0.6866 \\
S+D+R+F+H & 0.7397/0.6797 & 0.7423/0.6734 & 0.7543/0.6864 & 0.7617/0.6921 & 0.7680/0.6994 \\
S+D+P+F+H & 0.7134/0.6538 & 0.7123/0.6523 & 0.7267/0.6690 & 0.7397/0.6723 & 0.7503/0.6862 \\
S+R+P+F+H & 0.7179/0.6553 & 0.7170/0.6496 & 0.7295/0.6693 & 0.7442/0.6817 & 0.7522/0.6862 \\
D+R+P+F+H & 0.6708/0.6257 & 0.6879/0.6370 & 0.6926/0.6424 & 0.7024/0.6510 & 0.7102/0.6577 \\
S+D+R+P+F+H & 0.7264/0.6647 & 0.7277/0.6595 & 0.7403/0.6770 & 0.7542/0.6854 & 0.7641/0.6972 \\
\end{longtable}
\endgroup


Missingness-only and median-only diagnostics accompany the full models.
Sparse P measurements never define a complete-case subset for every family.
At ten seconds P is explicitly unavailable, so it cannot be counted as an
observed structural phenomenon. A grouped bootstrap assesses the predefined
S+D+R versus S+D+R+F, S+D+R+H and S+D+R+F+H comparisons; its intervals are
conditional on the already trained models and observed source collection,
not guarantees for future generators.

\subsection{Exact-60-second recurrence association: descriptive only}
The completed long-context analysis uses 1,604 development recordings (1,208 Human and 396 Suno), five globally blocked group folds and seven F/H/M combinations. The 438 historical locked and 50 pilot recordings are excluded from fitting. The only development AI source is Suno: generator holdout is ineligible, source-transfer $J$ remains undefined, and no cross-generator conclusion can be drawn.

\begin{table}[htbp]
\centering\footnotesize\setlength{\tabcolsep}{4pt}
\caption{Exact60 descriptive group-CV results: each cell is source-pair-averaged AUC/BA, not source-held-out $J$. Training caps are global groups per source. These numbers must not be compared with 30s source-transfer AUC as a causal duration or recurrence gain.}
\label{tab:fhm60-group-only}
\begin{tabular}{lrrrrr}
\toprule Combination & 25 & 50 & 100 & 200 & All \\
\midrule
F & 0.8676/0.7685 & 0.8675/0.7692 & 0.8685/0.7749 & 0.8748/0.7900 & 0.8741/0.7882 \\
H & 0.6914/0.6475 & 0.7074/0.6767 & 0.7169/0.6808 & 0.7151/0.6812 & 0.7131/0.6646 \\
M & 0.5966/0.5648 & 0.6239/0.5863 & 0.6323/0.6173 & 0.6327/0.6078 & 0.6319/0.6133 \\
F+H & 0.8607/0.7836 & 0.8811/0.8070 & 0.8866/0.8252 & 0.8910/0.8191 & 0.8904/0.8180 \\
F+M & 0.8736/0.7928 & 0.8875/0.8112 & 0.8887/0.8146 & 0.8938/0.8119 & 0.8925/0.8121 \\
H+M & 0.7012/0.6580 & 0.7262/0.6707 & 0.7328/0.6933 & 0.7331/0.6881 & 0.7315/0.6782 \\
F+H+M & 0.8630/0.7957 & 0.8930/0.8225 & 0.8986/0.8314 & 0.9014/0.8357 & 0.9014/0.8349 \\
\bottomrule
\end{tabular}
\end{table}

These are within-observed-source associations. Recurrence may reflect repertoire, loops, instrumentation or production in addition to generation. Old S/D/R/P features were not recomputed over this exact60 cohort, so this run does not isolate an M increment over the old pipeline. Missing H/M observations are retained under training-fold imputation and missingness diagnostics, not removed through complete-case filtering.


The separate 10-second S/D/R/F/H classification run is still in progress at
this checkpoint. Its 10,141-item extraction is complete, but no partial score
is promoted to a final result here. P has no eligible ten-second observation.

\subsection{Measurement admission outcomes and artifact references}

B uses a third-party VocalSet breath annotation layer, not assumed ground
truth for physiology. Silent/uncertain labels are not positive inhalations.
The singer-heldout measurement benchmark below is separate from detector use. A
requires independently validated singing phone timing and meaningful note
context; speech alignment confidence or word error rate is not a substitute.
T requires genuine isolated sources and pitch/register controls; a Demucs
\texttt{other} stem is not one instrument. These gates remain explicit and are
not silently replaced by spectral noise or flux summaries.

The B acquisition independently verifies 244 original VocalSet members and
their annotation/PCM/CRC/SHA evidence, but not the full archive MD5. Only 114
clips have a named reviewer and positive review time. A filename duplicated in
the ZIP has two different waveforms; the duration-matching member conflicts
with its singer directory and is conservatively excluded before fitting.
The resulting prospective cohort has 113 clips, 20 singer groups, 141 merged
high/medium-confidence events and 38 zero-event clips. Model-prefilled,
single-reviewer annotations and only three songs limit independence. The
provisional B admission gate is evaluated independently of AI/Human classification.
\begin{table}[htbp]
\centering\small
\caption{Completed external B benchmark: leave-one-singer-out, 113 clips, 20 singers and 141 merged reference events. This is event measurement, not AI-music classification.}
\label{tab:breath-external-failed}
\begin{tabular}{lr}
\toprule Metric & Value \\
\midrule
Event precision & 0.098278 \\
Event recall & 0.971631 \\
Event F1 & 0.178502 \\
Frame balanced accuracy & 0.912017 \\
Hard-negative frame false-positive fraction & 0.336281 \\
TP / FP / FN events & 137 / 1,257 / 4 \\
Zero-event clips with at least one false event & 38 / 38 \\
False events in zero-event clips & 352 \\
\bottomrule
\end{tabular}
\end{table}

The predeclared gate requires event precision and F1 both at least 0.70. It fails decisively despite high recall and frame BA. Every zero-event clip has a false alarm; 3,042 of 9,046 labeled hard-negative frames are positive. High frame-level discrimination is therefore not sufficient for reliable breath-event organization. No post-hoc threshold or event-duration rule is tuned on this benchmark, and B is not admitted to the AI/Human classifier. The audit verifies all 230 artifact hashes and event/model records; it does not validate event measurement.


An A data-admission review distinguishes NUS-48E phone-boundary annotations
from CSD note/syllable annotations. Neither inspected package alone validates
the full phone--note coordination chain. No A measurement or attribution result
is claimed, no access agreement was accepted on the user's behalf, and A is
not included in the classification lattice.

An initial T implementation was tested on 52 controlled pairs from 153 audited
reference-source files, rendering 208 six-second stimuli with nonoverlapping
three-second source chunks. At the known switch boundary, switched-source
distance exceeds same-source distance for 44/52 pairs (median difference
0.2423; descriptive pair-bootstrap 95\% interval [0.1877, 0.2981]); gain/polarity changes
are zero. However, the whole-clip p90 trajectory difference is positive for
only 32/52 pairs, with interval [0, 0.1225]. A known-boundary sensitivity result
does not establish a reliable unlocalized clip detector. Vocal-track identity
is not verified singer identity, and approximate register matching is not pitch
invariance. T is consequently not promoted as a validated natural-identity
feature in the F/H classifier lattice.
The 52 pair rows reuse 25 of 78 source identities across roles (up to three
appearances). Thus the pair-row interval is conditional/descriptive, not an
independent-source confidence interval; future inference needs source-disjoint
pairs or dependency-aware resampling.

The local extension root is
\filepath{/Users/yi/Documents/code/music/artifacts/audio_phenomena_expansion_20260907/}.
Paths in the following table resolve against it. Full waveforms retain exact
remote paths in the input/item manifests. Every numerical conclusion should be
read with the cited coverage and status records.

\begin{longtable}{>{\raggedright\arraybackslash}p{24mm}>{\raggedright\arraybackslash}p{61mm}>{\raggedright\arraybackslash}p{64mm}}
\caption{Extension code, outputs and numerical-analysis references.}\\
\toprule Responsibility & Code / protocol & Output / analysis \\
\midrule\endfirsthead
\toprule Responsibility & Code / protocol & Output / analysis \\
\midrule\endhead
Exact inputs & \filepath{code/audio_inputs.py}; \filepath{code/extract_features.py}
& \filepath{results/features_10s_v1/}; \filepath{results/features_30s_v1/}; \filepath{audit/extraction_30s_v1.json} \\
Phase & \filepath{code/phase_features.py}; \filepath{code/external_phase_validation.py}
& \filepath{results/external_phase_validation/REPORT.md}; \filepath{pair_metrics.csv} within that directory \\
Tonal path / recurrence & \filepath{code/musical_features.py}; \filepath{code/external_musical_validation.py}
& \filepath{results/external_musical_validation_v1/}; \filepath{manifests/long_cohort/summary_60s.json} \\
Pitch tracking & \filepath{code/benchmark_pitch_tracking.py}; \filepath{PITCH_TRACKING_BENCHMARK_EN.md}
& \filepath{results/external_pitch_tracking/overall_metrics.json}; aligned contours, per-clip and per-singer metrics within that directory \\
Pitch microstructure & \filepath{code/pitch_microstructure_features.py}; \filepath{V_MICROSTRUCTURE_EN.md}
& \filepath{results/external_pitch_microstructure_v1/overall_metrics.json}; failed observability gate and all contour comparisons \\
Timbral continuity & \filepath{code/timbre_features.py}; \filepath{code/external_timbre_validation.py}
& \filepath{results/external_timbre_validation_v1/summary.json}; known-boundary sensitivity and weak whole-clip results \\
All combinations & \filepath{code/evaluate_new_phenomena.py}; \filepath{preregistration/fh_30s_v2/}
& \filepath{results/evaluation_fh_30s_v2/}; source-pair metrics, per-record predictions, folds and frozen models \\
Result presentation & \filepath{code/summarize_fh_evaluation.py}
& \filepath{results/presentation_fh_30s_v2/all_63x5_numeric_results.csv}; \filepath{RESULTS_FH_EN.md} within that directory \\
Uncertainty & \filepath{code/bootstrap_matched_deltas.py}
& \filepath{results/uncertainty_fh_30s_v2/}; conditional paired global-group bootstrap \\
Source sensitivity & \filepath{code/source_resampling_sensitivity.py}
& \filepath{results/source_sensitivity_fh_30s_v1/}; post-hoc crossed-source composition sensitivity, no refit \\
Independent audit & \filepath{code/audit_fh_results.py}
& \filepath{AUDIT_FH_RESULTS_EN.md}; \filepath{audit/fh_30s_v2_audit_report.json}; 16,984 checks, zero errors \\
Breath admission & \filepath{code/benchmark_breath_events_v2.py}; \filepath{BREATH_BENCHMARK_PROTOCOL_EN.md}
& \filepath{results/external_breath_events_v2/summary.json}; \filepath{BREATH_BENCHMARK_RESULTS_EN.md}; failed gate and independent 230-artifact audit \\
Long recurrence & \filepath{code/evaluate_new_phenomena_v3.py}; \filepath{preregistration/fhm_60s_v3/}
& \filepath{results/evaluation_fhm_60s_v3/}; \filepath{AUDIT_FHM_60S_RESULTS_EN.md}; seven-by-five descriptive grid, J undefined \\
Articulation admission & \filepath{A_ADMISSION_REVIEW_EN.md}
& Prospective annotation requirements and unverified phone--note chain; no classifier result \\
Pending gates & \filepath{BREATH_BENCHMARK_PROTOCOL_EN.md}; \filepath{EXPERIMENT_PROTOCOL.md}
& \filepath{CURRENT_STATE.md}; \filepath{RUN_LOG.md}; external validation and exclusion ledgers \\
\bottomrule
\end{longtable}

The preregistration's generated filename retains the word \texttt{draft}, but
the operative v2 receipt contains \texttt{status=frozen} and the reviewed
contract hash. A stale remote-base mismatch was detected before scoring; its
unfrozen v1 receipt is retained as invalid history. The corrected run imports
the byte-identical archived authoritative evaluator, not the stale remote copy.
